% =====================================================================
%  整理者　macanine <macanine@163.com>
%  仓库　　github.com/macanine/zju-linear-algebra-papers
%  来源　　解析据 papers/2020-2021-秋冬-期中-甲.tex 撰写（原卷见 sources/2016-2025-期中-甲-历年真题汇编.pdf 第 8 页）
% =====================================================================

\begin{solution}
记 $T_m$ 为主对角全为 $4$、次对角全为 $-1$、上对角全为 $-3$ 的 $m$ 阶三对角行列式（约定 $T_0=1$）。按第一行展开：元素 $-3$ 的余子式 $M_{12}=-T_{m-2}$（消去首行与第二列后呈分块上三角，左上角为 $-1$），故
\[
T_m=4T_{m-1}+(-3)(-1)M_{12}=4T_{m-1}-3T_{m-2}\qquad(m\ge2),
\]
且 $T_1=4$，$T_2=4\times4-3=13$。特征方程 $r^{2}-4r+3=0$ 的根为 $r=1,3$，设 $T_m=\alpha3^{m}+\beta$，由 $T_0=1$、$T_1=4$ 得 $\alpha+\beta=1,\ 3\alpha+\beta=4$，解出 $\alpha=\dfrac32$、$\beta=-\dfrac12$，即
\[
T_m=\frac{3^{m+1}-1}{2}.
\]

再算题面所印的 $D_n$。把它的 $(n-1,n)$ 元记作 $p$：原卷此处印作 $p=3$，而其余各行的上对角元均为 $-3$。按最后一行展开（末行只有 $(n,n-1)=-1$、$(n,n)=4$ 非零）：
\[
D_n=(-1)^{n+(n-1)}(-1)M_{n,n-1}+(-1)^{2n}\cdot4\cdot M_{n,n}=M_{n,n-1}+4M_{n,n}\qquad(n\ge2).
\]
其中删去末行、末列后剩下的恰是 $T_{n-1}$，即 $M_{n,n}=T_{n-1}$；删去末行、第 $n-1$ 列后，前 $n-2$ 行前 $n-2$ 列构成 $T_{n-2}$，第 $1$ 至 $n-2$ 行在第 $n$ 列均为零，只有第 $n-1$ 行末元为 $p$，故该子式呈分块下三角，$M_{n,n-1}=p\,T_{n-2}$。于是
\[
D_n=4T_{n-1}+p\,T_{n-2}.
\]
若上对角元统一为 $-3$（即 $p=-3$，这才是题的本意），则 $D_n=4T_{n-1}-3T_{n-2}=T_n=\dfrac{3^{n+1}-1}{2}$，正是题面要证的结论。但若按原卷印出的 $p=3$ 逐字计算，则
\[
D_n=4T_{n-1}+3T_{n-2}=2(3^{n}-1)+\frac{3(3^{n-1}-1)}{2}=\frac{5\cdot3^{n}-7}{2}\qquad(n\ge2),
\]
与题面结论不符（$n=3$ 时按印出的矩阵算得 $64$，题面结论为 $40$；$n=2$ 时印出的矩阵为 $\begin{pmatrix}4&3\\-1&4\end{pmatrix}$，算得 $19\ne13$）。
\ans{上对角元全为 $-3$ 时 $D_n=\dfrac{3^{n+1}-1}{2}$，即题面所要证的结论；严格按原卷印出的 $3$ 计算，则 $D_1=4$，$D_2=19$（$n=2$ 时唯一的超对角元恰在 $(n-1,n)$ 处、取 $3$），$n\ge2$ 时 $D_n=4T_{n-1}+3T_{n-2}=\dfrac{5\cdot3^{n}-7}{2}$，该闭式对一切 $n\ge1$ 成立。}
\rem{原卷第 1 题的上对角元在 $(n-1,n)$ 处印作 $3$，与其余各行的 $-3$ 不对称，是原卷笔误；按“忠实转录不订正”的约定题干保留原样，作答时须点明，否则照原样代入得不到题面结论。}
\end{solution}

\begin{solution}
先说明 $B$ 可逆。按第二行展开
\[
|B|=(-1)\cdot\begin{vmatrix}1&1\\2&3\end{vmatrix}=-1\ne0 .
\]
由 $2BA+CB=O$ 得 $2BA=-CB$，两边左乘 $B^{-1}$，得 $2A=-B^{-1}CB$，即
\[
A=-\frac12B^{-1}CB .
\]
于是
\[
A+3E=B^{-1}\Bigl(3E-\frac12C\Bigr)B ,
\]
两端取行列式，由相似矩阵行列式相等得
\[
|A+3E|=\Bigl|3E-\frac12C\Bigr| .
\]
$C=\operatorname{diag}(1,2,-1)$，故 $3E-\dfrac12C=\operatorname{diag}\Bigl(\dfrac52,2,\dfrac72\Bigr)$，从而
\[
|A+3E|=\frac52\times2\times\frac72=\frac{35}{2}.
\]
\ans{$|A+3E|=\dfrac{35}{2}$。}
\end{solution}

\begin{solution}
$BA=O$ 表示 $B$ 的每一行 $y^{T}$ 都满足 $y^{T}A=0$，即 $B$ 的各行（看作列向量）都属于 $N(A^{T})$。

\textbf{先由“存在非零 $B$”定 $t$。}若 $|A|\ne0$，则 $A$ 可逆，由 $BA=O$ 两边右乘 $A^{-1}$ 得 $B=O$，与 $B$ 非零矛盾。故必须 $|A|=0$。按第一行展开
\[
|A|=1\cdot(36-t^{2})-2\cdot(18-3t)+3\cdot(2t-12)=-t^{2}+12t-36=-(t-6)^{2},
\]
故 $|A|=0$ 给出 $t=6$。

\textbf{再定 $r(B)$。}$t=6$ 时 $A$ 的三行都成比例（均为 $(1,2,3)$ 的倍数），故 $r(A)=1$。由秩－零化度定理
\[
\dim N(A^{T})=3-r(A)=2 .
\]
$B$ 的各行都在 $N(A^{T})$ 中，故 $r(B)\le2$；又 $B\ne O$，故 $r(B)\ge1$。这里 $N(A^{T})$ 恰是方程 $y_1+2y_2+3y_3=0$ 的解空间，取一行非零解（如 $(-2,1,0)^{T}$）作 $B$ 的唯一非零行，得 $r(B)=1$；取两个线性无关解 $(-2,1,0)^{T}$、$(-3,0,1)^{T}$ 作 $B$ 的前两行、第三行取零，得 $r(B)=2$。两种秩都确实能取到。
\ans{$t=6$；$r(B)$ 的可能取值为 $1$ 或 $2$。}
\end{solution}

\begin{solution}
关键是先看出 $(1,1)$ 这个解。由 $a+b+c=0$，
\[
a\cdot1+b\cdot1+c=0,\qquad b\cdot1+c\cdot1+a=0,
\]
故 $(x_1,x_2)=(1,1)$ 满足前两个方程；而前两个方程的系数行列式
\[
\begin{vmatrix}a&b\\b&c\end{vmatrix}=ac-b^{2}\ne0,
\]
由 Cramer 法则，前两个方程有且只有这一个解。

因此整个方程组至多有一解，且有解当且仅当 $(1,1)$ 也满足第三式。注意 $(1,1)$ 满足方程 $px_1+qx_2+r=0$ 当且仅当 $p+q+r=0$：前两式的系数之和分别为 $a+b+c$ 与 $b+c+a$，正是它们都等于 $0$ 才使 $(1,1)$ 成为解。第三式的三个系数也应是 $a,b,c$ 的一个排列，即循环形式 $cx_1+ax_2+b=0$（三系数之和仍为 $a+b+c=0$），代入 $(1,1)$ 得 $c+a+b=0$，自动成立。于是方程组有唯一解 $(1,1)$，与 $a,b,c$ 的具体取值无关。

不过要指出：原卷第三式印作 $cx_1+bx_2+b=0$，其系数之和为 $c+2b=b-a$（用了 $a+b+c=0$），一般不为零。按这一印法，当 $a\ne b$ 时 $(1,1)$ 不满足第三式，方程组无解（例如 $a=1,b=0,c=-1$，第三式成为 $-x_1=0$，与前两式给出的 $x_1=1$ 矛盾）；只有 $a=b$（此时 $c=-2a$，由 $b^{2}\ne ac$ 得 $a\ne0$）时方程组有唯一解，解仍为 $(1,1)$。
\ans{方程组的唯一解为 $(x_1,x_2)=(1,1)$。（第三式取循环形式 $cx_1+ax_2+b=0$ 时，题面的“有唯一解”对一切满足 $b^{2}\ne ac$、$a+b+c=0$ 的 $a,b,c$ 成立；若按原卷印出的 $cx_1+bx_2+b=0$，则仅当 $a=b$ 时有解，解同样是 $(1,1)$。）}
\rem{原卷第 4 题第三式印作 $cx_1+bx_2+b=0$，系数之和不等于 $a+b+c$，照原样只在 $a=b$ 时方程组有解；题干按约定保留原样，此处按“三式系数取 $a,b,c$ 的排列”的本意作答并点明差异。}
\end{solution}

\begin{solution}
先由 $A^{*}$ 读出 $|A|$。$A^{*}$ 是下三角矩阵，故 $|A^{*}|=1\cdot1\cdot1\cdot8=8$。此处 $n=4$，伴随矩阵的行列式满足 $|A^{*}|=|A|^{n-1}=|A|^{3}$，故 $|A|^{3}=8$，$|A|=2\ne0$，即 $A$ 可逆，且由 $AA^{*}=|A|E$ 得
\[
A^{-1}=\frac{1}{|A|}A^{*}=\frac12A^{*}
=\begin{pmatrix}
\frac12&0&0&0\\
0&\frac12&0&0\\
\frac12&0&\frac12&0\\
0&-\frac32&0&4
\end{pmatrix}.
\]

由 $AB=B+3A$ 得 $(A-E)B=3A$。两边左乘 $A^{-1}$，注意 $A$ 与 $B$ 不可交换、只能整体左乘：
\[
(E-A^{-1})B=3E .
\]
又 $E-A^{-1}=\dfrac12(2E-A^{*})$，而
\[
2E-A^{*}=\begin{pmatrix}1&0&0&0\\0&1&0&0\\-1&0&1&0\\0&3&0&-6\end{pmatrix}
\]
仍是下三角矩阵，$|2E-A^{*}|=-6\ne0$，故 $E-A^{-1}$ 可逆。解 $(2E-A^{*})X=E$ 得
\[
(2E-A^{*})^{-1}=\begin{pmatrix}1&0&0&0\\0&1&0&0\\1&0&1&0\\0&\frac12&0&-\frac16\end{pmatrix},
\]
于是
\[
B=3(E-A^{-1})^{-1}=6(2E-A^{*})^{-1}
=\begin{pmatrix}6&0&0&0\\0&6&0&0\\6&0&6&0\\0&3&0&-1\end{pmatrix}.
\]
验算：把上式代回，$(A-E)B=3A$，即 $AB=B+3A$，说明结论正确。
\ans{$B=\begin{pmatrix}6&0&0&0\\0&6&0&0\\6&0&6&0\\0&3&0&-1\end{pmatrix}$。}
\end{solution}

\begin{solution}
\textbf{(1)} $A$ 的每一列都等于 $e=(1,1,\dots,1)^{T}$，即 $A=ee^{T}$，而 $e^{T}e=n$，故
\[
A^{2}=(ee^{T})(ee^{T})=e\,(e^{T}e)\,e^{T}=n\,ee^{T}=nA .
\]
于是对 $f(x)=x^{2}+ax$ 有 $f(A)=A^{2}+aA=(n+a)A$。因 $A\ne O$，故 $f(A)=O$ 当且仅当 $a=-n$，即 $f(x)=x^{2}-nx$。（$f(A)$ 与 $A$ 同为 $n$ 阶方阵，题面“二阶零方阵”应作“$n$ 阶零方阵”理解，即 $O$；若按 $2\times2$ 字面理解，则 $n\ne2$ 时不可能。）

\textbf{(2)} 由 $A^{2}=nA$ 归纳证明 $A^{k}=n^{k-1}A$：$k=1$ 时显然；设对 $k$ 成立，则
\[
A^{k+1}=A^{k}A=n^{k-1}A^{2}=n^{k-1}\cdot nA=n^{k}A .
\]
取 $k=100$ 得 $A^{100}=n^{99}A$。

\textbf{(3)} $A^{3}=A^{2}A=nA\cdot A=nA^{2}=n^{2}A$；又 $E$ 与 $A$ 可交换，直接展开：
\[
(A+E)^{3}=A^{3}+3A^{2}+3A+E=n^{2}A+3nA+3A+E=E+(n^{2}+3n+3)A .
\]

\textbf{(4)} 由 $A^{2}=nA$ 直接验算
\[
(A+E)\Bigl(E-\frac{1}{n+1}A\Bigr)
=E+A-\frac{1}{n+1}A-\frac{1}{n+1}A^{2}
=E+A-\frac{1}{n+1}A-\frac{n}{n+1}A=E,
\]
故 $A+E$ 可逆，且
\[
(A+E)^{-1}=E-\frac{1}{n+1}A .
\]
\ans{(1) $f(x)=x^{2}-nx$（$a=-n$）；(2) $A^{100}=n^{99}A$；(3) $(A+E)^{3}=E+(n^{2}+3n+3)A$；(4) $(A+E)^{-1}=E-\dfrac{1}{n+1}A$。}
\end{solution}

\begin{solution}
\textbf{充分性。}设存在 $C$ 使 $A=ABC$。乘一个矩阵不会增大秩，故
\[
r(A)=r(ABC)\le r(AB)\le r(A),
\]
两端夹逼得 $r(A)=r(AB)$。

\textbf{必要性。}设 $r(A)=r(AB)$。$AB$ 的每一列都是 $A$ 的列的线性组合（$AB$ 的第 $j$ 列等于 $A$ 乘以 $B$ 的第 $j$ 列），故
\[
\operatorname{Col}(AB)\subseteq\operatorname{Col}(A);
\]
而 $\dim\operatorname{Col}(AB)=r(AB)=r(A)=\dim\operatorname{Col}(A)$，于是 $\operatorname{Col}(AB)=\operatorname{Col}(A)$。

因此 $A$ 的每一列 $a_j$ 都是 $AB$ 的列的线性组合，即存在列向量 $c_j$ 使 $a_j=AB\,c_j$（$j=1,2,\dots,n$）。令 $C=(c_1,c_2,\dots,c_n)$，则 $C$ 为 $n$ 阶实方阵，且
\[
A=(a_1,\dots,a_n)=(ABc_1,\dots,ABc_n)=AB\,C .
\]
\qed
\end{solution}

\begin{solution}
取矩阵的等价标准形：存在可逆矩阵 $U,V$ 与对角矩阵
\[
D=\operatorname{diag}(d_1,\dots,d_r,0,\dots,0)\qquad(d_1d_2\cdots d_r\ne0,\ r=r(A))
\]
使得 $UAV=D$，即 $A=U^{-1}DV^{-1}$。在中插入 $U^{-T}U^{T}=E$：
\[
A=U^{-1}DU^{-T}\cdot U^{T}V^{-1}.
\]
令 $S=U^{-1}DU^{-T}$，$P=U^{T}V^{-1}$。由 $U,V$ 可逆知 $P$ 可逆；又 $D$ 是对角矩阵，
\[
S^{T}=(U^{-1}DU^{-T})^{T}=U^{-1}D^{T}U^{-T}=U^{-1}DU^{-T}=S,
\]
即 $S$ 对称。于是 $A=SP$，其中 $S$ 对称、$P$ 可逆（$r=0$ 时 $D=O$，$S=O$ 同样对称）。
\qed
\end{solution}
